% ============================================================
% UGA-Themed Beamer Template (Clean MWE)
% ============================================================

\documentclass[aspectratio=169,professionalfonts]{beamer}

% ---------- COLORS (swap with official UGA hex if you like) ----------
\usepackage{xcolor}
\definecolor{UGARed}{HTML}{BA0C2F}   % Verify with brand guide
\definecolor{UGABlack}{HTML}{000000}
\definecolor{UGAGray}{HTML}{8A8D8F}

% ---------- THEME & PROGRESS DOTS ----------
\usetheme{default}
\useinnertheme{circles}
\useoutertheme[subsection=false]{miniframes}
\usefonttheme{professionalfonts}
\setbeamertemplate{navigation symbols}{}
% Head/foot colors
\setbeamercolor{subsection in head/foot}{fg=white,bg=UGARed}
\setbeamercolor{headline}{fg=white, bg=UGARed}

\setbeamercolor{section in head/foot}{fg=white,bg=UGARed}
\setbeamercolor{section in head/foot shaded}{fg=white,bg=UGARed}
\setbeamercolor{footline}{bg=UGABlack,fg=white}
\setbeamercolor{frametitle}{fg=UGABlack}
\setbeamercolor{title}{fg=UGABlack}
\setbeamercolor{normal text}{fg=UGABlack}
\setbeamercolor{structure}{fg=UGARed}

% --- HEADLINE: slide dots grouped by section ---
\setbeamersize{mini frame size=0.16cm,mini frame offset=0pt}
% Show section labels above the clickable slide dots.
\setbeamertemplate{section in head/foot}[default]
\setbeamertemplate{section in head/foot shaded}[default][75]
\setbeamertemplate{mini frame}{%
  \tikz[baseline=-0.5ex] \filldraw[draw=white,fill=white,line width=0.6pt] (0,0) circle (2pt);%
}
\setbeamertemplate{mini frame in current subsection}{%
  \tikz[baseline=-0.5ex] \filldraw[draw=white,fill=UGARed,line width=0.6pt] (0,0) circle (2pt);%
}
\setbeamertemplate{mini frame in other subsection}{%
  \usebeamertemplate{mini frame in current subsection}%
}
\setbeamertemplate{mini frame in other section}{%
  \usebeamertemplate{mini frame in current subsection}%
}
\setbeamertemplate{headline}{%
  \begin{beamercolorbox}[wd=\paperwidth]{headline}
    \vskip1ex
    \insertnavigation{\paperwidth}
    \vskip3ex
  \end{beamercolorbox}
}

% Simple footline with slide numbers
\setbeamertemplate{footline}{
  \leavevmode%
  \hbox{%
    \begin{beamercolorbox}[wd=\paperwidth,ht=2.5ex,dp=1ex,leftskip=1em,rightskip=1em]{footline}%
      \insertshortauthor{} \hfill \insertframenumber/\inserttotalframenumber
    \end{beamercolorbox}}%
}

% ---------- PACKAGES ----------
\usepackage{amsmath,amssymb,mathtools}
\usepackage{graphicx}
\usepackage{booktabs}
\usepackage{bbm}
\usepackage{array}
\usepackage{adjustbox}
\usepackage{hyperref}
\usepackage{tikz}
\usetikzlibrary{arrows.meta, positioning}
\usepackage{tcolorbox}
\usepackage{eurosym}
\usepackage{threeparttable}
% ---------- METADATA (fill these later) ----------
\title{Optimal Dynamic Pricing With Preferences for Novelty}
\subtitle{}
\author{Tate Mason}
\institute{The University of Georgia}
\date{\today}

% ---------- MACROS ----------
% Numbered equation with label
\newcommand{\beq}[1]{\begin{equation}\label{#1}}
\newcommand{\eeq}{\end{equation}}
% Accent color highlight
\newcommand{\hl}[1]{\textbf{\textcolor{UGARed}{#1}}}
% Simple figure placeholder macro
\newcommand{\figph}[2][0.9\linewidth]{%
  \begin{figure}
    \centering
    \fbox{\rule{0pt}{0.55#1}\rule{#1}{0pt}}%
    \caption{#2}
  \end{figure}
}

% ============================================================

% Keep navigation links in the theme's white foreground on the red header.
\hypersetup{colorlinks=true, linkcolor=., urlcolor=UGARed}


% Short-version navigation only. Original frame content is preserved verbatim.
\newcommand{\shortbuttons}{}
\newcommand{\shortmaincount}{25}
\IfFileExists{assets/quantified_theory/plain/optimal_prices_by_gamma.pdf}{\renewcommand{\shortmaincount}{33}}{}
\newif\ifshortappendix
\shortappendixfalse
\setbeamertemplate{footline}{%
  \leavevmode\vbox{\hbox{%
  \begin{beamercolorbox}[wd=\paperwidth,ht=3.8ex,dp=.6ex,rightskip=1em]{normal text}%
    \hfill{\scriptsize\shortbuttons}
  \end{beamercolorbox}}%
  \hbox{\begin{beamercolorbox}[wd=\paperwidth,ht=2.5ex,dp=1ex,leftskip=1em,rightskip=1em]{footline}%
    \insertshortauthor{}\hfill
    \ifshortappendix Appendix\else\insertframenumber/\shortmaincount\fi
  \end{beamercolorbox}}}%
}


\begin{document}

% Original numbered slide 0
\gdef\shortbuttons{}
\begin{frame}[plain,noframenumbering,label=short-0]
  \titlepage
\end{frame}


\section{Introduction}

% Original numbered slide 1
\gdef\shortbuttons{}
\begin{frame}[label=short-1]{Motivation}
  \begin{itemize}
    \item Firms have more access than ever to consumer information
    \item Allows for targeted discounting based on preferences derived from repeat interactions
    \item Targeted discounts allow firms to steer consumers toward particular products.
  \end{itemize}
  \begin{tcolorbox}
    How do firms price based on repeated interactions with consumers and how are profits affected?
  \end{tcolorbox}
\end{frame}


% Original numbered slide 2
\gdef\shortbuttons{}
\begin{frame}[label=short-2]{Research Question}
  \begin{tcolorbox}
    How can firms price discriminate via consumer preferences for novelty and habit and how much is information about these preferences worth?
  \end{tcolorbox}
  How do I answer?
  \begin{itemize}
    \item Solve a theoretical model to understand firm's behavior
    \item Estimate consumer's preferences with Nielsen data, recovering the distribution of novelty/habit preferences
    \item Use the estimated distribution in the pricing model to understand how firm's revenues are impacted by this sort of targeting
  \end{itemize}
\end{frame}


% Original numbered slide 3
\gdef\shortbuttons{}
\begin{frame}[label=short-3]{Literature}
  \begin{columns}[T]
    \begin{column}{0.48\textwidth}
      \textbf{Targeting:}
      \begin{itemize}
        \item Dub\'e (2008)
        \item Dub\'e et. al (2022)
        \item Anand \& Shachar (2009)
        \item Karle \& Reisinger (2024)
        \item Durrmeyer, Fournel, Samano (2026)
      \end{itemize}
    \end{column}
    \begin{column}{0.48\textwidth}
      \textbf{Estimating preferences for novelty:}
      \begin{itemize}
        \item Dub\'e (2004)
        \item Dub\'e et. al (2021)
        \item Chintagunta (1999)
        \item Lattin (1987)
        \item McAllister (1982)
        \item Ackerberg (2003)
      \end{itemize}
    \end{column}
  \end{columns}
\end{frame}


\section{Motivation from Data}

% Original numbered slide 4
\gdef\shortbuttons{\hyperlink{short-5}{\beamergotobutton{Flavor Popularity}}}
\begin{frame}[label=short-4]{HHI}
  \centering
    \includegraphics[width=.95\textwidth,height=.74\textheight,keepaspectratio]{../writeup/1108_writeup/assets/results/summary_statistics/flavor_hhi_histogram.pdf}
\end{frame}


\section{Supply Theory}

% Original numbered slide 6
\gdef\shortbuttons{}
\begin{frame}[label=short-6]{Supply Setting}
  \begin{itemize}
    \item Monopolistic firm producing a single product $j$
    \item Firm is profit maximizing, setting prices according to that objective
    \item A consumer is born into the world with a type $r\in\{H,L\}$, determining how much utility they derive from switching and an initial state $h_0\in\{0,j\}$
    \item Consumers live for two periods, purchasing the inside good or opting for the outside good in each period
    \item The firm knows initial states and types, setting prices to individual consumers by type and history
  \end{itemize}
\end{frame}


% Original numbered slide 7
\gdef\shortbuttons{}
\begin{frame}[label=short-7]{Consumer's Problem}
  \begin{tcolorbox}[
    boxsep=0pt,
    left=2pt, right=2pt,
    top=1pt, bottom=1pt,
    before skip=2pt, after skip=2pt
    ]
    \setlength{\abovedisplayskip}{2pt}
    \setlength{\belowdisplayskip}{2pt}
    \begin{gather*}
      V^{inside}_{jt} = \underbrace{\beta}_{\text{Constant}} + \overbrace{\gamma_r\mathbbm{1}\{h=0\}}^{\text{Utilty from Switching}} - \underbrace{\alpha p_t(h)}_{\text{Price Disutility}}, \ V^{Outside}_{0t} = \underbrace{\gamma_r\mathbbm{1}\{h=j\}}_{\text{Utility from Switching}} \\
      v_t(x) = V_{jt}^{inside} - V_{0t}^{outside} = \beta + \underbrace{\gamma[\mathbbm{1}\{h=0\} - \mathbbm{1}\{h=j\}]}_{q(h)} - \alpha p_t(h)
    \end{gather*}
  \end{tcolorbox}
  \begin{itemize}\setlength{\itemsep}{2pt}
    \item Prices, set by the firm, and switching utility depend on history $h=x_{t-1}$
    \item Utility for the outside good is normalized to zero
  \end{itemize}
\end{frame}


% Original numbered slide 8
\gdef\shortbuttons{}
\begin{frame}[label=short-8]{Firm's Problem}
  \begin{tcolorbox}
    [
    boxsep=0pt,
    left=2pt, right=2pt,
    top=1pt, bottom=1pt,
    before skip=2pt, after skip=2pt
    ]
    \setlength{\abovedisplayskip}{2pt}
    \setlength{\belowdisplayskip}{2pt}
    \begin{gather*}
      \pi_t(h) = \max_{p_t(h)} s_t(h)[p_t(h) - c] \\
      \Pi(h) = \underbrace{s_1(h)(p_1(h) - c)}_{\text{Margin in Period 1}} + \underbrace{\delta\{s_1(h)\pi_2(j) + (1-s_1(h))\pi_2(0)\}}_\text{{Continuation Value}}
    \end{gather*}
  \end{tcolorbox}
  \begin{itemize}
    \item Firm profits are comprised of period one margins, given history of consumption, and the value of that consumer in the different states in period two
    \item The firm will set prices according to consumer type and history to maximize expected profit
  \end{itemize}
\end{frame}


% Original numbered slide 10
\gdef\shortbuttons{\hyperlink{short-9}{\beamergotobutton{Terminal FOC}}\hspace{.35em}\hyperlink{short-11}{\beamergotobutton{Comparative statics}}\hspace{.35em}\hyperlink{short-52}{\beamergotobutton{Elasticities}}}
\begin{frame}[label=short-10]{First Period Optimization}
  \begin{tcolorbox}
    [
    boxsep=0pt,
    left=2pt, right=2pt,
    top=1pt, bottom=1pt,
    before skip=2pt, after skip=2pt
    ]
    \setlength{\abovedisplayskip}{2pt}
    \setlength{\belowdisplayskip}{2pt}
    \begin{align*}
      \frac{d\Pi}{d p_1(h_0)} &= s_1 + \frac{\partial}{\partial p_1(h)}s_1(h_0)[p_1(h_0) - c + \delta(\pi_2(j) - \pi_2(0))] \\
      p_1^*(h_0) &= c + \frac{1}{\alpha(1-s_1^*(h_0))} - \delta(\pi_2(j) - \pi_2(0)) \\
      \Pi^*(h) &= \delta\pi_2(0) + \frac{s_1^*(h_0)}{\alpha(1-s_1^*(h_0))}
    \end{align*}
  \end{tcolorbox}
\end{frame}


% Original numbered slide 14
\gdef\shortbuttons{\hyperlink{short-12}{\beamergotobutton{Specification}}\hspace{.35em}\hyperlink{short-13}{\beamergotobutton{Numerical table}}}
\begin{frame}[label=short-14]{Optimal Prices}
  \centering
  \includegraphics[width=0.85\textwidth]{assets/theory_delta1/optimal_prices_by_gamma.pdf}
\end{frame}


% Theory plots regenerated at delta=1.
\gdef\shortbuttons{\hyperlink{short-14}{\beamergotobutton{Optimal prices}}\hspace{.35em}\hyperlink{short-13}{\beamergotobutton{Numerical table}}}
\begin{frame}[label=short-prices-profit]{Optimal Price Ratios and Profit}
  \centering
  \begin{columns}[T,onlytextwidth]
    \begin{column}{.49\textwidth}
      \includegraphics[width=\linewidth,height=.66\textheight,keepaspectratio]{assets/theory_delta1/all_state_price_ratios.pdf}
    \end{column}
    \begin{column}{.49\textwidth}
      \includegraphics[width=\linewidth,height=.66\textheight,keepaspectratio]{assets/theory_delta1/all_state_realized_profit.pdf}
    \end{column}
  \end{columns}
  \vspace{.3em}
  {\footnotesize $\delta=1$: no discounting. Flavor is $j$; outside is everything except that flavor.}
\end{frame}


% Original numbered slide 18
\gdef\shortbuttons{}
\begin{frame}[label=short-18]{Two Pricing Regimes: The Same Consumer Types}
  \small
  \begin{columns}[T,onlytextwidth]
    \begin{column}{.48\textwidth}
      \small\textbf{Household-specific: type observed}
      \begin{itemize}
        \item Observe type $\gamma_r$ and initial history $h_0$.
        \item Choose $p_{1r}(h_0)$ and then $p_{2r}(h_1)$.
      \end{itemize}
    \end{column}
    \begin{column}{.48\textwidth}
      \small\textbf{Common: type hidden}
      \begin{itemize}
        \item Observe history and know type weights $\mu_r$.
        \item One offer across types with the same history.
        \item Fixed beliefs restrict learning from purchases.
      \end{itemize}
    \end{column}
  \end{columns}
  \medskip
  \begin{tcolorbox}
    Consumers keep their type in both periods. Evaluate both policies on these
    same persistent types; the gain is household-specific minus common revenue.
  \end{tcolorbox}
  {\footnotesize $\gamma_r\in\{-1.25,0,1.25\}$; equal weights $1/3$ or permutations of $(.50,.25,.25)$.
  In all cases, $h_0$ and subsequent purchases are observed; $\delta=1$, $c=0$.}
\end{frame}


% Original numbered slide 20
\gdef\shortbuttons{\hyperlink{short-19}{\beamergotobutton{Profit levels}}\hspace{.35em}\hyperlink{short-53}{\beamergotobutton{Price paths}}\hspace{.35em}\hyperlink{short-54}{\beamergotobutton{Revenue evaluation}}}
\begin{frame}[label=short-20]{Theory: Revenue Gain from Observing Type}
  \centering
  \includegraphics[width=.98\textwidth,height=.66\textheight,keepaspectratio]{assets/theory_delta1/gamma_uncertainty/information_value.pdf}
  \par\smallskip\footnotesize Expected revenue per consumer, evaluated under the same persistent types.
  Initially outside means outside the focal flavor; it is an initial state, not a population share.
\end{frame}


\section{Data}

% Original numbered slide 21
\gdef\shortbuttons{\hyperlink{short-22}{\beamergotobutton{Activity}}}
\begin{frame}[label=short-21]{Data Sources}
  \textbf{Nielsen HomeScan}
  \begin{itemize}
    \item Provides household level purchase data
    \item Covers years 2006-2015
    \item Products chosen, quantity chosen, timing of trips
  \end{itemize}
  \textbf{Nielsen RetailScan}
  \begin{itemize}
    \item Provides store level product data
    \item For specified product category: prices, markets, product attributes
  \end{itemize}
  \textbf{Sample Description}
  \begin{itemize}
    \item Must be single-person household
    \item Must be seen at least 2 times
    \item Purchases must be made at Nielsen store (prices available)
    \item Flavor attribute must be present
  \end{itemize}
\end{frame}


% Original numbered slide 23
\gdef\shortbuttons{}
\begin{frame}[label=short-23]{Sample Statistics: Demographics}
  \centering
  \begin{adjustbox}{max totalsize={\textwidth}{.83\textheight},center}
    \begin{minipage}{\textwidth}
      \input{../writeup/1108_writeup/assets/results/summary_statistics/full_panel_vs_yogurt.tex}
    \end{minipage}
  \end{adjustbox}
\end{frame}

% Original numbered slide 24
\gdef\shortbuttons{}
\begin{frame}[label=short-24]{Switching Stats}
    \centering
    \begin{adjustbox}{max totalsize={\textwidth}{.75\textheight},center}
    \begin{minipage}{\textwidth}
      \input{../RESULTS_ONLY/summary_statistics/aggregated_flavors/yogurt_dynamics.tex}
    \end{minipage}%
    \end{adjustbox}
\end{frame}


% Original numbered slide 25
\gdef\shortbuttons{}
\begin{frame}[label=short-25]{Persistence Across Aggregated Flavor Groups}
  \centering
  \includegraphics[height=0.69\textheight]{../writeup/1108_writeup/assets/results/summary_statistics/aggregated_flavors/flavor_transitions.pdf}
  \par\small Plain repeats in approximately 80\% of eligible transitions, versus 53\% for berry.
\end{frame}


\section{Demand Estimation}

% Original numbered slide 26
\gdef\shortbuttons{}
\begin{frame}[label=short-26]{Estimate Heterogeneity in Price, Novelty and Habit}
  \begin{itemize}
    \item Weekly choices: other yogurt, berry, plain, or nonpurchase.
    \item Households differ in price sensitivity and responses to recent and
    longer purchase histories.
    \item Flavor intercepts are common across households and differ across flavors.
  \end{itemize}
  \begin{tcolorbox}
    Recover the distribution of switching preferences, then use it in the
    two-period pricing model.
  \end{tcolorbox}
\end{frame}


% Original numbered slide 27
\gdef\shortbuttons{\hyperlink{short-average-history}{\beamergotobutton{Average-history specification}}}
\begin{frame}[label=short-27]{Demand Model}
  \footnotesize
  \[
    U_{ijt}=\beta_0+\beta_jX_{jt}+\gamma_iN_{ijt}+\phi_iH_{ijt}
             +\alpha_i p_{jt}+\varepsilon_{ijt},
    \qquad U_{i0t}=\varepsilon_{i0t}.
  \]

  Household $i$'s share of yogurt purchases in flavor $j$ in week $t$ is
  \[
    w_{ijt}=\frac{y_{ijt}}{\sum_k y_{ikt}}.
  \]

  \begin{columns}[T,onlytextwidth]
    \begin{column}{.58\textwidth}
      \textbf{Novelty}
      \[
        N_{ijt}=\begin{cases}
          1-w_{ij,t-1}, & \text{if yogurt bought in $t-1$},\\
          0, & \text{otherwise}.
        \end{cases}
      \]
    \end{column}
    \begin{column}{.35\textwidth}
      \textbf{Habit: fitted specification}
      \[
        H_{ijt}=\sum_{s<t}w_{ijs}.
      \]
    \end{column}
  \end{columns}
\end{frame}

\gdef\shortbuttons{}
\begin{frame}[label=short-identification]{Separating Recent and Longer Purchase History}
  \begin{itemize}
    \item \textbf{Novelty:} changes in last week's basket.
    \item \textbf{Habit:} differences in earlier purchase histories, even with the same last-week basket.
    \item \textbf{Simplifying assumption:} common flavor tastes; heterogeneous novelty, habit, and price sensitivity.
  \end{itemize}
\end{frame}
\note{Common flavor intercepts keep the specification tractable and focus on differences in responses to purchase history. Persistent household flavor tastes could explain some observed repetition; interpretation of history coefficients is conditional on abstracting from that heterogeneity and on assumptions about utility shocks. The uncontrolled fit assumes prices are exogenous to omitted utility shocks. Histories start at zero at the first observed purchase, leaving unobserved pre-sample history as a limitation.}


% Original numbered slide 29
\gdef\shortbuttons{\hyperlink{short-likelihood}{\beamergotobutton{Likelihood}}\hspace{.35em}\hyperlink{short-demand-fit}{\beamergotobutton{Fit summary}}\hspace{.35em}\hyperlink{short-28}{\beamergotobutton{Probabilities}}\hspace{.35em}\hyperlink{short-price-distribution}{\beamergotobutton{Price sensitivity}}}
\begin{frame}[label=short-29]{Estimation Choices}
  \small
  \textbf{Purchase occasions}
  \begin{itemize}
    \item Count packages within observed yogurt-purchase weeks; calendar weeks without purchases are not added.
    \item Impute outside choices as $\max(7-Q_{it},0)$, so each retained week contributes $\max(7,Q_{it})$ choices.
  \end{itemize}
  \medskip
  \textbf{Persistent household heterogeneity}
  \[
    \gamma_i\sim\mathcal N(\mu_\gamma,\sigma_\gamma^2),\quad
    \phi_i\sim\mathcal N(\mu_\phi,\sigma_\phi^2),\quad
    \alpha_i=-\exp(\mu_{\log a}+\sigma_{\log a}\nu_{\alpha i}).
  \]
  \begin{itemize}
    \item Independent coefficient draws, held fixed across weeks; $\nu_{\alpha i}\sim\mathcal N(0,1)$.
    \item Simulated maximum likelihood integrates over household types using 128 draws.
  \end{itemize}
\end{frame}


% Original numbered slide 32
\gdef\shortbuttons{\hyperlink{short-30}{\beamergotobutton{Estimates}}\hspace{.35em}\hyperlink{short-33}{\beamergotobutton{Discussion}}}
\begin{frame}[label=short-32]{Distributions}
  \begin{center}
    \includegraphics[width=.96\textwidth,height=.56\textheight,keepaspectratio]{assets/demand_1111_three/history_distributions.pdf}
  \end{center}
\end{frame}


% Original numbered slide 31
\gdef\shortbuttons{\hyperlink{short-wtp-interpretation}{\beamergotobutton{WTP interpretation}}}
\begin{frame}[label=short-31]{Implied Willingness to Pay}
  \centering
  \input{assets/demand_1111_three/implied_wtp.tex}
  \begin{itemize}
    \item Novelty WTP refers to a one-unit increase in $N$, holding other utility terms fixed; it is not the total compensation for a switch.
    \item Habit WTP refers to a one-unit increase in $H$. Population SD measures heterogeneity, not estimation uncertainty.
  \end{itemize}
\end{frame}


\section{Pricing Exercise}

% Original numbered slide 34
\gdef\shortbuttons{\hyperlink{short-theory-mapping}{\beamergotobutton{Parameter mapping}}}
\begin{frame}[label=short-34]{Quantifying the Theory with Estimated Preferences}
  \small
  \[
    v_t=\beta_f+\alpha p_t+\gamma\big[\mathbbm{1}\{h_t=0\}-\mathbbm{1}\{h_t=j\}\big].
  \]
  \begin{itemize}
    \item Plain versus not plain: use the estimated novelty distribution, mean price coefficient and plain intercept; hold habit at zero.
    \item Informed prices condition on $\gamma$; common prices use its distribution with fixed beliefs. Both adapt terminal prices to purchase history.
    \item Evaluate both policies on the same persistent types over two periods.
  \end{itemize}
  \medskip
  This is an illustrative quantification exercise.
\end{frame}

\gdef\shortbuttons{}


% Revised empirical information comparison; populated by aggregate-only refresh.
\IfFileExists{assets/quantified_theory/plain/optimal_prices_by_gamma.pdf}{%
\gdef\shortbuttons{}
\begin{frame}[label=short-revised-prices]{Quantified Theory: Optimal Prices across Novelty Types}
  \centering
  \includegraphics[width=.98\textwidth,height=.68\textheight,keepaspectratio]{assets/quantified_theory/plain/optimal_prices_by_gamma.pdf}
  \par\smallskip\footnotesize Same model and layout as supply, using estimated preference parameters.
  There are exactly two terminal prices, after plain and after not plain.
\end{frame}

\gdef\shortbuttons{}


\gdef\shortbuttons{\hyperlink{short-berry-paths}{\beamergotobutton{Berry}}\hspace{.35em}\hyperlink{short-other-paths}{\beamergotobutton{Other}}}
\begin{frame}[label=short-revised-paths]{Quantified Theory: Informed and Common Price Paths}
  \centering
  \includegraphics[width=.98\textwidth,height=.68\textheight,keepaspectratio]{assets/quantified_theory/plain/information_price_paths_comparison.pdf}
  \par\smallskip\footnotesize Solid paths average informed prices within novelty groups; dashed paths show common prices.
  Both regimes choose the terminal price after observing the first-period binary choice.
\end{frame}

\gdef\shortbuttons{}


\gdef\shortbuttons{\hyperlink{short-all-flavors}{\beamergotobutton{All flavors}}}


\gdef\shortbuttons{\hyperlink{short-revised-table}{\beamergotobutton{Revenue table}}\hspace{.35em}\hyperlink{short-prior-details}{\beamergotobutton{Distribution weights}}}
\begin{frame}[label=short-revised-gain]{Knowing Novelty Adds More Revenue for Initially Outside Consumers}
  \centering
  \includegraphics[width=.98\textwidth,height=.68\textheight,keepaspectratio]{assets/quantified_theory/plain/information_value.pdf}
  \par\smallskip\footnotesize Informed minus common revenue, evaluated on the same persistent types.
  Prior sensitivities use one novelty group with probability one-half.
\end{frame}

\gdef\shortbuttons{}

}{%
\gdef\shortbuttons{}
\begin{frame}[label=short-revised-pending]{Binary Plain-versus-Not Pricing: Rerun Pending}
  \begin{itemize}
    \item Saved estimates are inputs to the binary model, including known habit.
    \item Recompute household-specific and common prices over two binary purchase periods.
    \item Only novelty is hidden; price sensitivity and habit are known.
    \item Report price changes and revenue gains in the same figure layouts as supply.
  \end{itemize}
  \begin{tcolorbox}
    Revised results require a pricing-only server run.
    Earlier full-basket results do not apply to this binary specification.
  \end{tcolorbox}
\end{frame}
}





% Original numbered slide 46
\gdef\shortbuttons{}
\begin{frame}[label=short-46]{Calibrated Plain Illustration: Revenue Levels}
  \centering
  \includegraphics[width=.98\textwidth,height=.67\textheight,keepaspectratio]{assets/aligned_pricing/information_friction_profits.pdf}
  \par\smallskip\footnotesize Dollars per model consumer over two periods.
  Both policies are evaluated under the same persistent types; $\delta=1$.
\end{frame}





% Original numbered slide 47


% Original numbered slide 57






\section{Conclusion}

% Original numbered slide 50
\gdef\shortbuttons{}
\begin{frame}[label=short-50]{What Does Knowing Novelty Change?}
  \begin{itemize}
    \item Estimated novelty preferences are negative on average and heterogeneous; average habit preference is positive.
    \item In the fitted-distribution illustration, information adds about 3.6\% for initially inside and 12.5\% for initially outside consumers.
    \item These magnitudes depend on the binary parameter mapping, exogenous-price assumption, and fixed-belief benchmark.
  \end{itemize}
  \begin{tcolorbox}
    Preferences for novelty and habit create opportunities for price discrimination.
  \end{tcolorbox}
\end{frame}


% Original numbered slide 51



\appendix
\input{appendix_current.tex}

\end{document}
